\appendix

\chapter{Installation and Repository}

The full source code for this project is publicly available at:

\begin{center}
\texttt{https://github.com/fionnmcgoldrick123/programming-quiz-app}
\end{center}

\section*{System Requirements}

The platform is containerised and requires the following software
to be installed on the host machine:

\begin{itemize}
    \item Docker (version 24.0 or later)
    \item Docker Compose (version 2.20 or later)
    \item Git
\end{itemize}

\section*{Quick Start}

\begin{enumerate}
    \item \textbf{Clone the repository:}
    \begin{lstlisting}[language=bash]
git clone https://github.com/fionnmcgoldrick123/programming-quiz-app
cd programming-quiz-app
    \end{lstlisting}

    \item \textbf{Create the environment configuration file:}
    \begin{lstlisting}[language=bash]
cp .env.example .env
    \end{lstlisting}

    \item \textbf{Set the following required variables inside \texttt{.env}:}
    \begin{itemize}
        \item \texttt{OPENAI\_API\_KEY} --- OpenAI API key for
        GPT-4o-mini generation
        \item \texttt{JWT\_SECRET} --- secret key for JWT signing
        (generate with \texttt{openssl rand -hex 32})
        \item \texttt{POSTGRES\_PASSWORD} --- database password
        \item \texttt{SMTP\_HOST}, \texttt{SMTP\_USERNAME},
        \texttt{SMTP\_PASSWORD}, \texttt{SMTP\_FROM} --- email
        server credentials for user verification
    \end{itemize}

    \item \textbf{Build and start all services:}
    \begin{lstlisting}[language=bash]
docker-compose up --build
    \end{lstlisting}

    The build process takes approximately two to three minutes on
    first run. On completion, the services are accessible at the
    following addresses:
    \begin{itemize}
        \item Frontend: \texttt{http://localhost:3000}
        \item Backend API: \texttt{http://localhost:8000}
        \item API documentation: \texttt{http://localhost:8000/docs}
    \end{itemize}

    \item \textbf{To stop all services:}
    \begin{lstlisting}[language=bash]
docker-compose down
    \end{lstlisting}
\end{enumerate}

\section*{Retraining the Machine Learning Models}

The serialised model artefacts are included in the repository and
loaded automatically at startup. Should retraining be required, the
difficulty classifier training script is located at
\texttt{ml\_models/difficulty\_classifier/difficulty\_predictor.py}
and requires the \texttt{CODENET\_PATH} environment variable to point
to the root of the IBM Project CodeNet dataset. The topic tag
classifier training script is located at
\texttt{ml\_models/tag\_classifier/train.py} and reads from the
included \texttt{leetcode.csv} file. Both scripts save their output
to \texttt{difficulty\_model.pkl} and \texttt{topic\_model.pkl}
respectively.